\documentclass[tikz, border=10pt]{standalone}
\usetikzlibrary{calc, decorations.pathmorphing, arrows.meta, patterns}

\begin{document}

% =========================================================
% ГОЛОВНІ ПАРАМЕТРИ КЕРУВАННЯ ТА КОЛЬОРИ
% =========================================================
\def\Phase{140}        % Фаза коливань у градусах (0..360)
\def\MaxLines{5}     % Максимальна кількість ліній поля / зарядів
\def\Amp{1.2}        % Амплітуда зміщення маятника (см)

% Палітра кольорів
\colorlet{ColorE}{red}     % Електричне поле E
\colorlet{ColorB}{blue!40}       % Магнітне поле B
\colorlet{ColorCoil}{cyan!80!black}    % Котушка індуктивності L
\colorlet{ColorMass}{gray} % Вантаж m
\colorlet{ColorSwitch}{red}            % Перемикач
\colorlet{ColorVelocity}{red!80!black} % Вектор швидкості v
% =========================================================

\begin{tikzpicture}[>=stealth, font=\large]

    % --- Фізичні обчислення (автоматичні) ---
    \pgfmathsetmacro{\xnorm}{cos(\Phase)}       % Нормоване зміщення x/x0 ∈ [-1, 1]
    \pgfmathsetmacro{\inorm}{-sin(\Phase)}      % Нормований струм I/I0 ∈ [-1, 1]

    % Кількість ліній полів та зарядів
    \pgfmathsetmacro{\numE}{int(round(abs(\xnorm) * \MaxLines))}
    \pgfmathsetmacro{\numB}{int(round(abs(\inorm) * \MaxLines))}

    % Фізичне зміщення маси в см та геометрія осі x
    \pgfmathsetmacro{\xShift}{\xnorm * \Amp}
    \pgfmathsetmacro{\xEquil}{2.2}
    \pgfmathsetmacro{\xPos}{\xEquil + \xShift}

    % =========================================================
    % 1. ЕЛЕКТРИЧНИЙ КОНТУР (LC)
    % =========================================================

    \def\PlateH{0.12} % Товщина пластин

    % --- ВЕРХНЯ ПЛАСТИНА ---
    % Передня грань (з урахуванням товщини)
    \filldraw[fill=gray!25, draw=black, thick, line join=round]
        (0, 1.2-\PlateH) -- (2.0, 1.2-\PlateH) -- (2.0, 1.2) -- (0, 1.2) -- cycle;
    % Права грань
    \filldraw[fill=gray!35, draw=black, thick, line join=round]
        (2.0, 1.2-\PlateH) -- (2.6, 1.6-\PlateH) -- (2.6, 1.6) -- (2.0, 1.2) -- cycle;
    % Верхня грань
    \filldraw[fill=gray!10, draw=black, thick, line join=round]
        (0, 1.2) -- (2.0, 1.2) -- (2.6, 1.6) -- (0.6, 1.6) -- cycle;

    % --- НИЖНЯ ПЛАСТИНА ---
    % Передня грань
    \filldraw[fill=gray!25, draw=black, thick, line join=round]
        (0, 0.0-\PlateH) -- (2.0, 0.0-\PlateH) -- (2.0, 0.0) -- (0, 0.0) -- cycle;
    % Права грань
    \filldraw[fill=gray!35, draw=black, thick, line join=round]
        (2.0, 0.0-\PlateH) -- (2.6, 0.4-\PlateH) -- (2.6, 0.4) -- (2.0, 0.0) -- cycle;
    % Верхня грань
    \filldraw[fill=gray!10, draw=black, thick, line join=round]
        (0, 0.0) -- (2.0, 0.0) -- (2.6, 0.4) -- (0.6, 0.4) -- cycle;

    \node[above left] at (0.6, 1.6) {$C$};

    \node[above left] at (0.6, 1.6) {$C$};

    % --- Заряди та Електричне поле E ---
    \ifnum \numE > 0
        \foreach \i in {1,...,\numE} {
            \pgfmathsetmacro{\t}{\i / (\numE + 1)}
            \pgfmathsetmacro{\px}{0.4 + \t * 1.8}

            \ifdim \xShift pt > 0pt
                % Позитивний заряд зверху, поле вниз
                \node[red, font=\small] at (\px, 1.35) {$+$};
                \node[blue, font=\small] at (\px, 0.25) {$-$};
                \draw[->, thick, ColorE] (\px, 1.15) -- (\px, 0.45);
            \else
                % Негативний заряд зверху, поле вгору
                \node[black, font=\small] at (\px, 1.35) {$-$};
                \node[black, font=\small] at (\px, 0.25) {$+$};
                \draw[->, thick, ColorE] (\px, 0.45) -- (\px, 1.);
            \fi
        }
        \node[right, ColorE, font=\bfseries] at (2.4, 0.8) {$\Efield$};
    \fi

    % Провідники контуру
    \draw[thick] (1.3, 1.4) -- (1.3, 2.3) -- (4.2, 2.3);
    \draw[thick] (1.3, -\PlateH) -- (1.3, -0.9) -- (2.0, -0.9);
    \draw[thick] (2.8, -0.9) -- (4.2, -0.9);

    % Перемикач
%    \draw[thick, ColorSwitch] (2.0, -0.9) circle (2pt);
%    \draw[thick, ColorSwitch] (2.8, -0.9) circle (2pt);
%    \ifnum \numB = 0
%        \draw[thick, ColorSwitch] (2.0, -0.9) -- (2.7, -0.6);
%    \else
%        \draw[thick, ColorSwitch] (2.0, -0.9) -- (2.8, -0.9);
%    \fi

    % Котушка індуктивності L
    \draw[thick, ColorCoil, decoration={coil, aspect=0.55, segment length=3.2mm, amplitude=3.5mm}, decorate]
        (4.2, 2.3) -- (4.2, -0.9);
    \node[right=8pt] at (4.2, 0.7) {$L$};

    % --- Магнітне поле B у котушці ---
    \ifnum \numB > 0
        \foreach \j in {1,...,\numB} {
            \pgfmathsetmacro{\offset}{(\j - (\numB+1)/2) * 0.15}
            \pgfmathsetmacro{\bendAngle}{\offset * 25}

            \ifdim \inorm pt > 0pt
                % Знизу вгору
                \draw[->, thick, ColorB] (4.2+\offset, -1.5) to[bend left=\bendAngle] (4.2+\offset, 3.0);
            \else
                % Зверху вниз (використовуємо bend left для збереження форми дуги)
                \draw[->, thick, ColorB] (4.2+\offset, 3.0) to[bend right=\bendAngle] (4.2+\offset, -1.5);
            \fi
        }
        \node[left=4pt, ColorB, font=\bfseries] at (3.9, 0.7) {$\Bfield$};
    \fi

    % Динамічні підписи струму I та заряду q
    \ifnum \numB = 0
        \node[above] at (2.75, 2.3) {$I = 0$};
    \else
        \node[above] at (2.75, 2.3) {$I \neq 0$};
    \fi

    \ifdim \xShift pt > 0pt
        \node[above right] at (2.3, 1.5) {$+q_0$};
        \node[below right] at (2.3, 0.1) {$-q_0$};
    \else\ifdim \xShift pt < 0pt
        \node[above right] at (2.3, 1.5) {$-q_0$};
        \node[below right] at (2.3, 0.1) {$+q_0$};
    \fi\fi

    % =========================================================
    % 2. ПРУЖИННИЙ МАЯТНИК
    % =========================================================
    \begin{scope}[shift={(6.2,0)}]
        % Стіна
        \fill[pattern=north east lines] (0,-0.2) rectangle (-0.25, 1.8);
        \draw[thick] (0,-0.2) -- (0, 1.8);

        % Центр вантажу
        \coordinate (MassCenter) at (\xPos, 0.5);

        % Пружина
        \draw[thick, decoration={coil, aspect=0.4, segment length=2.2mm, amplitude=2.8mm}, decorate]
            (0, 0.5) -- (\xPos - 0.5, 0.5);
%        \node[above=6pt] at (1.0, 0.5) {$k$};

        % Маса (Вантаж m)
        \draw[thick, fill=ColorMass, rounded corners=1pt]
            (\xPos - 0.5, 0.1) rectangle (\xPos + 0.5, 0.9);
        \node at (MassCenter) {$m$};

        % Швидкість v
        \ifnum \numB = 0
            \node[above=15pt] at (MassCenter) {$v = 0$};
        \else
            \pgfmathsetmacro{\vDir}{sign(\inorm)}
            \draw[->, very thick, ColorVelocity]
                (\xPos, 1.0) -- (\xPos + \vDir*0.7, 1.0)
                node[midway, above] {$v$};
        \fi

        % Вісь координат x
        \draw[->, thick] (0, -0.4) -- (4.2, -0.4) node[below] {$x$};
        \draw[thick] (\xEquil, -0.3) -- (\xEquil, -0.5) node[below=2pt] {\scriptsize $0$};

        % Розмірна лінія зміщення x0
        \draw[gray, dashed] (\xEquil, -0.3) -- (\xEquil, -1.1);
        \draw[gray, dashed] (\xPos, -0.3) -- (\xPos, -1.1);

        \ifdim \xShift pt = 0pt
        \else
            \draw[<->, >=Latex] (\xEquil, -0.95) -- (\xPos, -0.95)
                node[midway, above=-2pt, font=\small] {$x$};
        \fi

    \end{scope}

\end{tikzpicture}
\end{document}